\section{Bifurcation of multisections in every degree}\label{sec:walls}

Fix a closed Riemann surface $C$ of genus $g\ge2$ and a smooth complex
plane bundle $N$ of degree $2g-2$. We work with a fixed Dolbeault
operator $\bar\partial_N$ on $N$. A smooth fiber map
$B_t(x,v)\in\Lambda_x^{0,1}C\otimes N_x$, defined near the zero
section with $B_t(x,0)=0$, determines an almost complex structure with
holomorphic projection to $C$ and complex affine fibers. For an \'etale
cover $f:Y\to C$, the graph of $s\in\Gamma(f^*N)$ is holomorphic
precisely when
\begin{equation}\label{eq:wall-section}
  \bar\partial_{f^*N}s+f^*B_t(\cdot,s)=0.
\end{equation}
Here pullback in the form factor uses $df$. In a holomorphic frame,
the horizontal $(0,1)$ lift of $\partial_{\bar z}$ has vertical
$(1,0)$ component $-B_{t,\bar z}$. This gives the asserted structure
and equation. Over a compact interval these structures are tamed on
a common tube about the zero section. Indeed the differential of the
fiberwise radial primitive, extended by a unitary horizontal
distribution, restricts to the standard symplectic form on each fiber.
Adding a sufficiently large multiple of base area gives tameness on
the zero section and hence on a tube. The normal operator of the core is
$D_t=\bar\partial_N+d_vB_t(\cdot,0)$. Every smooth path of real
Cauchy--Riemann operators on $N$ with this symbol arises in this way.

Take as $D_0,D_1$ two of the spin models of Section~\ref{sec:spin}
with core signs $+1$, and with difference characters $0$ and
$w\ne0$ respectively. One can arrange positive parity for a fixed
nonzero $w$ by twisting the two spin lines simultaneously by a class
pairing nontrivially with $w$ if necessary. The two smooth complex
bundles are isomorphic since they have equal rank and degree. Fix a
complex isomorphism with $N$ for each. Let $a_1$ be a simple
nonseparating circle representing the Poincar\'e dual of $w$
modulo two. Complete it to surface generators $a_i,b_i$ with
relation $\prod_i[a_i,b_i]=1$. For each $m\ge2$ let
$f_m:Y_m\to C$ be the connected \'etale cover on $m$ letters defined by
\begin{equation}\label{eq:wall-quotient}
 a_1\longmapsto(1\ 2),\qquad a_2\longmapsto(1\ 2\ \cdots\ m),
 \qquad b_i,a_j\longmapsto 1\quad(j\ge3\text{ for }a_j).
\end{equation}
These permutations generate $S_m$, and $Y_m$ has genus
$h_m=1+m(g-1)$. Write $d_m=|\operatorname{Aut}_C(f_m)|$, so
$d_2=2$ and $d_m=1$ for $m\ge3$.

\begin{theorem}\label{thm:walls}
For every smooth path from $D_0$ to $D_1$, the fiberwise linear graph
equations have primitive holomorphic multisections with projection
$f_m$ arbitrarily close to the zero section at some parameter, for
every $m\ge2$. Here primitive means that the graph map is simple and
its image projects to $C$ with degree $m$.

For a residual choice of path of normal operators, each $f_m$ has an
odd finite number of transverse crossings in its trace-free summand.
These crossings occur at different parameters for different $m$ and
the base operator is invertible at each one. Fix such a path of first
jets. For a residual choice of the higher fiber jets in
\eqref{eq:wall-section}, each crossing has a branch of primitive
multisections. In a neighborhood of the multiple cover the change,
from the left to the right, in their signed count is
\begin{equation}\label{eq:wall-orientation}
  \frac{\sigma(f_m)^- - \sigma(f_m)^+}{d_m},
\end{equation}
where $\sigma(f_m)^\pm$ are the determinant signs of $f_m^*D_t$ on
the adjacent sides. Multiplying each change by the sign of the base
operator at the crossing and summing along the path gives $2/d_m$.

One may instead choose the nonlinear terms in the space of fiberwise
odd maps. A residual choice in this space gives a birth or a death at
every one of the stated crossings. There are two nearby simple
supports on the occupied side for $m\ge3$ and one for $m=2$.
\end{theorem}

The count in \eqref{eq:wall-orientation} concerns the specified
projection cover at source genus $h_m$. Indeed the nonconstant part
of an arbitrary map of projection degree $m$ has arithmetic genus at
least $h_m$ by Riemann--Hurwitz. Equality requires a single \'etale
component. A contracted genus-zero tree attached to it has an unstable
leaf, and a cycle in the dual graph would increase the genus.
The topological type of the resulting \'etale projection is discrete.
For a single elementary wall in the space of compatible almost complex
structures on a closed symplectic Calabi--Yau threefold, Bai and
Swaminathan proved the local signed change $\pm2/|\operatorname{Aut}(f)|$,
including the standard symmetric-group wall. They also give unramified
dihedral examples in every odd degree \cite[Theorems~1.9 and~4.12,
Examples~4.3--4.4]{bai}. Here the endpoint
signs force walls of every degree for one path in the space of graph
equations. We prove the local assertion in that space and retain the
orientation of the change.

\subsection*{The trace-free summand}
Pushing the pullback normal operator down to $C$ amounts to twisting
$D_t$ by the real permutation local system of $f_m$. That local
system is $\RR\oplus W_m$, where $W_m$ is the sum-zero standard
representation of $S_m$ pulled back by \eqref{eq:wall-quotient}.
Both summands have index zero. This follows by deforming to a complex
Dolbeault operator and using that a complexified flat real local
system has degree zero. The complex orientations give
\begin{equation}\label{eq:wall-relative}
  \operatorname{sgn}(D_t^{W_m})=
  \operatorname{sgn}(f_m^*D_t)\operatorname{sgn}(D_t)
\end{equation}
when both operators on the right are invertible. The endpoint formula
of Section~\ref{sec:spin} gives $+1$ and $-1$ on the left, since the
monodromy of $a_1$ is odd. Thus some $D_t^{W_m}$ has a nonzero
kernel. For the linear graph equation any nonzero real multiple of
such a kernel element is a solution on $Y_m$ with zero fiberwise
trace.

A section on $Y_m$ pulls back to a section on the normal closure
invariant under $S_{m-1}$. If two distinct local sheets agree on an
open disk, the corresponding translates agree identically on the
normal closure. Indeed their difference satisfies a real linear
Cauchy--Riemann equation, obtained by integrating the fiber derivative
along the segment joining them in a common tube. Unique continuation
applies. The stabilizer would then properly contain $S_{m-1}$, which
is a maximal subgroup, so all sheets would agree identically. It
follows that every section which does not descend has pairwise distinct
values over a dense open set. Its graph map is simple and immersive
because its projection is \'etale. A nonzero trace-free section cannot
descend. Positive multiples tending to zero give distinct supports,
as can also be seen from their distinct maximum fiber norms.
The deck orders stated above follow from
$\operatorname{Aut}_C(f_m)=N_{S_m}(S_{m-1})/S_{m-1}$.

\subsection*{Choice of the path}
We use equivariant Brill--Noether transversality for real
Cauchy--Riemann operators \cite{wendl,dwbn}. Here are the precise
hypotheses for this application. The affine parameter space of normal
operators allows every real zeroth-order bundle map as an infinitesimal
variation, including variations supported in any given coordinate
disk, with arbitrary prescribed finite jet at an interior point.
The real Cauchy--Riemann symbol satisfies Wendl's condition
\cite[Theorem~2.5.3]{dwbn}. Off a set of infinite codimension the
operator then satisfies the equivariant Petri condition for any
specified finite list of inequivalent irreducible local systems
\cite[Proposition~1.6.6 and Theorem~1.6.17]{dwbn}. The operator and
its formal adjoint satisfy strong unique continuation by the
Cauchy--Riemann symbol and the zeroth-order similarity principle
(the adjoint has the conjugate symbol). In these statements
one may take the smooth Banach perturbations with prescribed jets
used for differential operators in \cite[Sections~2.2--2.3]{dwbn}.
Away from this exceptional set, the joint kernel stratum for the
trivial system and any finite selection of $W_m$ has real codimension
$\sum_i k_i^2$, where $k_i$ is the real nullity of the $i$th twist
\cite[Theorem~1.3.5 and Proposition~1.3.11]{dwbn}. Each standard representation is absolutely
irreducible and of real type. The different ranks distinguish the
$W_m$, and $W_2$ is distinguished from the trivial line by its
nontrivial monodromy.

Relative transversality for paths now gives transverse intersections
with the corank-one strata and avoidance of the higher-codimension
strata and of the exceptional set \cite[Proposition~1.B.5]{dwbn}. For clarity, variations of a path
supported in a short interior parameter interval realize the operator
variations just described at its center. Thus evaluation of paths
at an interior point is a submersion in these perturbation spaces.
One obtains the same residual assertion for smooth paths by the
usual local perturbation argument: for each finite list the required
property is open on a compact interval, by Fredholm reduction and
endpoint invertibility, and dense in each smooth neighborhood using
small smooth Banach perturbations there. To see openness one can
express the condition directly as simple transverse crossings with
no simultaneous nullities; this expression requires no condition on
jets away from the crossing set. Intersect the resulting open dense
sets over all finite initial lists. For each $W_m$ its crossing set
is closed in the compact interval and each of its points is isolated,
so it is finite. Its parity is odd by \eqref{eq:wall-relative}.
For any prescribed degree bound this gives only finitely many of the
selected crossings.

\subsection*{Nonlinear graphs and their signs}
Fix one of the crossings $t_*$ for a selected $f=f_m$. Let $\xi$ and
$\eta$ span the real kernel and adjoint kernel of $f^*D_{t_*}$ and
choose $L^2$ splittings. The crossing pairing
$\langle\eta, f^*\dot D_{t_*}\xi\rangle_{L^2}$ is nonzero.
Lyapunov--Schmidt reduction of \eqref{eq:wall-section} on
$W^{1,r}(Y_m,f^*N)$, $r>2$, solves the image equation uniquely for
the complementary part of $s=a\xi+z(t,a)$. Here $z(t,0)=0$ and
$z_a(t_*,0)=0$. Elliptic regularity gives smooth sections. The
remaining real equation is $aG(t,a)=0$. Its derivative $G_t(t_*,0)$
is the crossing pairing up to the nonzero normalization of the
cokernel coordinate. Thus all nearby nonzero solutions lie on a
smooth branch $t=T(a)$ through $T(0)=t_*$. Their nonzero kernel
projection rules out descent, since base sections are orthogonal to
$\xi$ in the fiberwise permutation splitting.

We verify that a higher fiber jet can give $T(a)-t_*$ a first nonzero
Taylor coefficient, keeping the normal path fixed. Put
$A=\dot D_{t_*}$, a real zeroth-order map. For some base point $x$
the sum $\sum_i u_i$ over its sheets is nonzero, where
$u_i=\langle\eta_i,A(x)\xi_i\rangle_{\RR}$ in the pulled-back
metrics. This follows by integration of the crossing pairing. Let
$r_1,\ldots,r_k$ be the distinct positive values of $|\xi_i|^2$ at
$x$. The matrix $(r_j^l)_{1\le l,j\le k}$ is invertible. Since
$u_i=0$ when $\xi_i=0$, there is an $l$ with $1\le l\le m$ such that
$\sum_i |\xi_i|^{2l}u_i\ne0$. Add to $B$ a small real multiple of
\begin{equation}\label{eq:wall-jet}
        \chi(t,x)|v|^{2l} A(x)v,
\end{equation}
where a nonnegative cutoff localizes the pairing to a disk on which
this sum keeps its strict sign, and equals a base cutoff at $t_*$.
This changes the coefficient of $a^{2l+1}$ in the reduced equation
at $t_*$ by a nonzero multiple of the perturbation, keeping lower
coefficients fixed. Indeed substitution of $s=a\xi+O(a^2)$ gives
the displayed pairing at that order. The induced change in $z$ starts
at the same order, where its linear image under $f^*D_{t_*}$ has
zero cokernel projection. If lower coefficients vanish, this changes
the coefficient of $a^{2l}$ in $T(a)-t_*$ by a nonzero multiple as
well, because $G_t(t_*,0)\ne0$. Finite order at most $2m$ is
therefore an open dense condition at this crossing, including in the
fiberwise odd subspace. The nonlinearities with a prescribed first
jet form an affine Fr\'echet space on a closed tube; intersect these
conditions over the countably many crossings. One may fix the germs
at both endpoints by using parameter cutoffs in \eqref{eq:wall-jet}.

Finite order implies that $a=0$ is the only solution at $t_*$ in a
small interval of kernel coordinates. The branch has no critical
point on either punctured half of a sufficiently small such interval.
On regular parameter slices to either side all solutions there are
isolated and regular. Brouwer degree of the reduced real map is
constant on these slices. The sign of the complementary invertible
block supplies the same orientation factor on both sides. The core
solution contributes $\sigma(f)^\pm$, so subtraction yields the
numerator in \eqref{eq:wall-orientation}. For \'etale projections the
tangential map $T_{Y_m}\to f^*T_C$ is a complex isomorphism and
cancels the source deformation factor with positive orientation. The
section sign is consequently the determinant sign for the graph as
a map with variable domain. The deck group acts freely on the
non-descending sections with stabilizer $S_{m-1}$ on the normal
closure, and its orbits parameterize their images. Deck changes
preserve the complex determinant orientation. Division by $d_m$
therefore gives the signed count of distinct simple curves.

In an odd family the splittings may be used unchanged under $s\mapsto
-s$. Uniqueness in the image equation gives $T(-a)=T(a)$. Its first
nonconstant term has even order, so one side has two solutions with
opposite kernel coordinates, and the other has none. In degree two
the nontrivial deck transformation acts by $-1$ on the kernel. Its
action exchanges the solutions, giving one simple support. In higher
degree the deck group is trivial and the two supports are distinct.
Finally at a selected crossing the sign of $D_t$ is constant.
Multiplying \eqref{eq:wall-orientation} by this sign gives the
left relative sign minus the right one divided by $d_m$, by
\eqref{eq:wall-relative}. The relative sign stays constant across
all other parameters at which $D_t^{W_m}$ is invertible, including
base-only crossings. The sum is the difference between its endpoint
values divided by $d_m$. This proves Theorem~\ref{thm:walls}.

\subsection*{A limiting operator on an infinite cover}
The infinitely many crossings have an analytic accumulation point even
though each individual projection cover has only finitely many. Give
$C$ a fixed Hermitian metric and give the pullback bundles and covers
their lifted metrics.

\begin{proposition}\label{prop:wall-limit}
For the generic paths of Theorem~\ref{thm:walls}, there is an interior
accumulation parameter $t_\infty$ of the selected crossing times
and a connected infinite \'etale cover of $C$ on which the pulled-back
$D_{t_\infty}$ has a nonzero bounded smooth solution. On this cover
the bottom of the spectrum of the Friedrichs extension of
$D_{t_\infty}^*D_{t_\infty}$ is zero.
\end{proposition}

\begin{proof}
At each spin endpoint write $D=\bar\partial+R$, with $R$ the
antilinear mass defined by the nondegenerate holomorphic symmetric
form $\gamma:N\to K_C N^*$ and the metric, as in
Section~\ref{sec:spin}. On any finite \'etale cover the integrated
real cross term of $\bar\partial s$ and $Rs$ vanishes: up to a fixed
area normalization it is the real part of the integral of
$\bar\partial(\gamma(s,s))$, where the canonical factor is pulled
through the covering differential. Hence
\[
        \|Ds\|_2^2=\|\bar\partial s\|_2^2+\|Rs\|_2^2
                       \ge c^2\|s\|_2^2
\]
for the positive least singular value bound $c$ of $R$ on $C$.
This bound is independent of the \'etale cover. An endpoint complex
bundle identification only changes the equivalent norms by constants
from the base. Along a fixed-symbol path the difference from the
endpoint is zeroth order and its uniform norm lifts unchanged. There
are therefore collars of both endpoints free of these crossings in
all degrees. Choose $m\to\infty$ and one crossing $t_m$ for each
$m$. A subsequence converges to an interior parameter $t_\infty$.

Take $\xi_m$ in the trace-free kernel on $Y_m$, normalized to have
maximum fiber norm one. Choose a point in the fiber of a fixed base
point within uniformly bounded distance of a maximum point, by lifting
a shortest path on $C$. The based covers have a subsequence converging
on bounded sets to an infinite cover, as can be seen directly from
\eqref{eq:wall-quotient}. Label the sheets cyclically, with the
chosen point labelled zero. The second generator acts by unit shift
and the first interchanges two consecutive labels. After passage to
a subsequence this interchange either occurs at fixed integer labels
or leaves every bounded set, in which case the first generator acts
trivially in the limit. In either case the limiting sheet set is
$\ZZ$ and the second generator acts by translation, so the cover is
connected and infinite. More explicitly, words of any bounded length
have eventually the same action at the distinguished sheet as on this
labelled limit. The resulting covering charts on each compact part
identify the lifted operator coefficients there, up to their smooth
dependence on $t_m$. The genus and the compact base are fixed.
Interior elliptic estimates on successively larger compact sets and a
diagonal subsequence give smooth local convergence of $\xi_m$ to a
solution $\xi$ with $|\xi|\le1$. The points of maximum remain in a
bounded set of the based covers, so $|\xi|$ attains one somewhere.

The limiting cover has area growth $O(R)$ on radius-$R$ balls. To see
this, use a fixed finite triangulation of $C$. Its translates at each
sheet have uniformly bounded area; after choosing a spanning tree,
crossing any edge moves the sheet by at most a fixed number of places
on $\ZZ$, since its monodromy is a fixed word in the surface
generators. A metric ball of radius $R$ reaches only $O(1+R)$ sheets:
lifting a fixed finite cover by evenly covered balls gives an edge
word of length $O(1+R)$ for any path in that ball of length at most
$R$. Let $\chi_R$ equal one on the radius-$R$ ball and vanish outside
the radius-$2R$ ball, with weak gradient bounded by $2/R$.
The Leibniz rule gives
\[
   \|D_{t_\infty}(\chi_R\xi)\|_2^2=O(R^{-1}),
     \qquad \liminf_{R\to\infty}\|\chi_R\xi\|_2^2>0.
\]
The cutoffs belong to the form domain by compact support and smooth
approximation. Their Rayleigh quotients tend to zero, proving the
spectral assertion. The construction gives boundedness of $\xi$;
its square integrability on the limiting cover is a separate question.
\end{proof}

There is also a spectral statement on the fixed base. On the Hilbert
direct sum over $m\ge2$ of the two $L^2$ spaces of the twist $W_m$
(sections and $(0,1)$-forms), take the self-adjoint direct sum of the
operators
\[
  \mathcal A_{m,t}=
  \begin{pmatrix}0&(D_t^{W_m})^*\\ D_t^{W_m}&0\end{pmatrix}.
\]
The direct sum has its natural graph domain. At a selected $t_m$ put
an $L^2$-unit kernel section into the section component of its own
summand. For a subsequence $t_m\to t_\infty$ these vectors are
orthonormal, and the norm of their image under the direct sum at
$t_\infty$ tends to zero: it is bounded by the uniform zeroth-order
norm of $D_{t_m}-D_{t_\infty}$ on $C$. Weyl's criterion puts zero in
the essential spectrum of this direct sum at $t_\infty$. At the two
endpoints the preceding uniform estimate applies to every summand.
Each twist has index zero, so its inverse and the adjoint inverse
have uniformly bounded norms. Thus the direct sum has a spectral gap
about zero at both endpoints.

\begin{corollary}\label{cor:torus}
For every $g\ge2$ a closed symplectic six-torus admits a smooth path
of tame almost-complex structures fixing an embedded genus-$g$ curve
in a primitive integral class $A$, with the bifurcations of
Theorem~\ref{thm:walls} in all the classes $mA$, $m\ge2$.
One may realize the fiberwise odd families and their signed changes.
\end{corollary}
\begin{proof}
Use the same smoothed symplectic surface and the same complex normal
identification for the whole path. All the graph structures agree as
complex structures along the zero section. Compactness of the parameter
interval gives a single tube on which they lie in the chosen chart
about the ambient compatible extension and are tame. Apply the cutoff
construction there, leaving them unchanged on a common smaller tube.
For every $m$ the multisections occur arbitrarily close to the zero
section, so they persist in this one compact realization. Their
integral homology classes are $mA$ by the tubular retraction.
\end{proof}

\subsection{Reduction to the normal operator}\label{sec:rescale}

Let $C$ be a closed embedded pseudoholomorphic curve of genus $g\ge2$
in an almost-complex six-manifold $(X,J)$ with $c_1(TX)[C]=0$.
Write $D$ for the normal Cauchy--Riemann operator on
$N=TX|_C/TC$. For an integer $M\ge1$, consider the condition
\begin{equation}\label{eq:normal-injectivity}
  \ker(f^*D)=0
  \quad\text{for every connected holomorphic cover }f:Y\to C
  \text{ of degree at most }M .
\end{equation}
Here $Y$ is smooth and closed, ramification is allowed, and the
identity is included. Write $J_D$ for the homogeneous almost-complex
structure on $N$ associated with $D$. Its projection to $C$ is
holomorphic, its fibers have their vector-space complex structures,
and a lift over $f$ is holomorphic exactly when its section of
$f^*N$ lies in $\ker(f^*D)$ \cite[Lemma~3.6]{dwcast}.

\begin{proposition}\label{prop:rescale}
Under \eqref{eq:normal-injectivity} there is a disk tube about $C$
in which the only cycle of degree $m$ is $m[C]$ for $1\le m\le M$.
Its point class $P_m(J)$ in $S$ equals $P_m(J_D)$.
The tube can be chosen uniformly over a smooth compact interval of
marked germs satisfying \eqref{eq:normal-injectivity} throughout,
and these classes are constant along the interval. If the condition
holds in every degree, the entire series of the germ has the same
unique lift respecting the cycle operations as that of its normal
model. This lift is constant under such deformations.
\end{proposition}

\begin{proof}
Use the marked surface to identify the cores of the family $J_t$.
Choose complex complements along the cores and identify the normal
bundles with a fixed smooth complex bundle $N$ over the interval.
In the resulting tubular charts put $s_r(x,v)=(x,rv)$ and
$J_{t,r}=s_r^*J_t$ for $0<r\le1$. In bundle coordinates the
tangent-to-normal block is $r^{-1}B_t(x,rv)$ with $B_t(x,0)=0$.
It extends smoothly in all variables as
\[
  \int_0^1 \partial_v B_t(x,urv)[v]\,du .
\]
The normal-to-tangent block is multiplied by $r$ and the diagonal
blocks are evaluated at $(x,rv)$. The limit $J_{t,0}$ is $J_{D_t}$
by the normal rescaling formula \cite[Lemma~3.10]{dwcast}.
For a change $F$ of tubular charts, $s_r^{-1}Fs_r$ extends at zero
to the induced map on the quotient normal bundles. Indeed, its base
component tends to the map on the cores and its vertical component
tends to the vertical derivative of $F$ applied to $v$. Thus the
identification is independent of the complement and is natural under
complex normal isomorphisms.

Realize a common disk of this compact family on the symplectic
six-torus of Proposition~\ref{prop:factors}, with $A=[C]$ primitive.
Indeed, in the symplectic splitting along the core the vertical
complex structure is fixed, and each induced complex structure on
$TC$ is compatible with its area form. The extension and cutoff in
that proposition therefore give a smooth family tamed by one form
$\omega$, agreeing with $J_{t,r}$ on a common closed disk of some
radius $2a>0$.

If there were no uniform isolating radius through degree $M$, there
would be $J_{t_i,r_i}$-holomorphic cycles $Z_i$ of degrees
$1\le m_i\le M$ with positive maximum fiber norms $b_i\to0$.
Pass to a subsequence with $m_i=m$ and $t_i\to t$. Put $u_i=b_i/a$.
Then $(s_{u_i}^{-1})_*Z_i$ is holomorphic for $J_{t_i,r_i u_i}$,
also when $r_i=0$ by homogeneity. It lies in the closed disk of
radius $a$ and meets its boundary. Dilation preserves the class
$mA$ by the tubular retraction, so its symplectic area is
$m\langle[\omega],A\rangle$. Geometric cycle compactness gives a
limit for $J_{t,0}$ in the same class, with maximum norm $a$ by
Hausdorff convergence \cite[Theorem~2.11 and Remark~2.13]{diw}.
This is the rescaling argument of \cite[\S6]{dwcast}.

Normalize an irreducible component of the limit off the zero
section. Its projection $f$ is nonconstant, since a complex vector
space contains no compact curve. All components thus have positive
projection degree, counted with their multiplicities to give $m$.
Our component gives a nonzero element of $\ker(f^*D_t)$ on a
smooth cover of degree at most $m$ \cite[Proposition~3.12]{dwcast}.
This contradicts \eqref{eq:normal-injectivity}. The argument uses
cycle compactness without imposing a bound on the genus.

In class $mA$ for $1\le m\le M$, containment of the support in a
smaller open disk selects just the section $m[C]$. It is open by Hausdorff
continuity and compact over the closed parameter rectangle. Its
characteristic class therefore transports along either parameter
by the interval relation \cite[\S2]{WueProduct}. Open excision
identifies the classes obtained using smaller tubes.
When \eqref{eq:normal-injectivity} holds for every $M$, the linear
model has only the zero section as a compact irreducible curve,
and each unramified pullback of $D$ is invertible by index zero.
Proposition~\ref{prop:factors} gives its balanced lift, and
Proposition~\ref{prop:full-transport} gives uniqueness among lifts
respecting the cycle operations. The coefficientwise equalities
just proved give the same lift for the original germ and for every
member of the interval.
\end{proof}
